% Canonical 6-slide Beamer fixture for the theme-authoring smoke test.
%
% Mirrors fixtures/canonical-deck-marp.md — same six slides exercise the
% same brand surfaces (h1/h2, bullets, section divider, code block, table,
% inverted closing).
%
% Build:
%   cp <generated>/beamer/beamerthemeHcw.sty <here>/
%   cp <generated>/beamer/hcw-logo.{png,pdf} <here>/  # raster sibling for pdfLaTeX
%   latexmk -pdf canonical-deck-beamer.tex
%
% The deck assumes the generated theme is placed next to this .tex.
\documentclass[11pt,aspectratio=169]{beamer}
\usetheme{Hcw}

\usepackage[utf8]{inputenc}
\usepackage[T1]{fontenc}
\usepackage{listings}
\usepackage{xcolor}

\title{Generated theme test}
\subtitle{Verifying brand projection}
\author{Andreas Drozd}
\date{14 May 2026}

\begin{document}

% ── Slide 1: title (lead variant) ─────────────────────────────────────────
\begin{frame}[plain,noframenumbering]
  \titlepage
\end{frame}

% ── Slide 2: bullets and emphasis ────────────────────────────────────────
\begin{frame}{Bullets and emphasis}
  The 6-slide canonical deck exercises every variant the generator emits:
  \begin{itemize}
    \item Title slide (lead variant)
    \item Section divider (invert variant)
    \item Bullet body
    \item Code block
    \item Table
    \item Inverted closing slide
  \end{itemize}
  Use \emph{italic} for emphasis, \textbf{bold} for terms, and \texttt{inline code} for identifiers.
\end{frame}

% ── Slide 3: section divider (invert variant) ────────────────────────────
% \section{} triggers the AtBeginSection divider page configured by the theme.
\section{Visual surface}

% ── Slide 4: code block ──────────────────────────────────────────────────
\begin{frame}[fragile]{Code block}
\begin{lstlisting}[language=,basicstyle=\ttfamily\small,backgroundcolor=\color{BrandBgAlt},frame=leftline,rulecolor=\color{BrandPrimary},framerule=2pt,xleftmargin=1em]
apiVersion: apps/v1
kind: Deployment
metadata:
  name: theme-authoring-test
spec:
  replicas: 3
  selector:
    matchLabels:
      app: test
\end{lstlisting}
  Note the brand-color inline \texttt{code} and the brand-color left rule on the block.
\end{frame}

% ── Slide 5: table ───────────────────────────────────────────────────────
\begin{frame}{Tables}
  \begin{tabular}{lll}
    \hline
    \textbf{\color{BrandPrimary}Token} & \textbf{\color{BrandPrimary}Value} & \textbf{\color{BrandPrimary}Where it lives} \\
    \hline
    primary       & \texttt{\#36566D} & h1, brand surfaces, code rule \\
    secondary     & \texttt{\#4B5059} & h2, body emphasis             \\
    primary\_soft & \texttt{\#7897AC} & h3, subtitles                 \\
    \hline
  \end{tabular}

  \vspace{1em}
  Header rows carry the brand-primary underline.
\end{frame}

% ── Slide 6: inverted closing slide ──────────────────────────────────────
% Beamer doesn't have a built-in "invert" body slide; we approximate with
% a beamercolorbox spanning the frame.
\begin{frame}[plain,noframenumbering]
  \begin{beamercolorbox}[wd=\paperwidth,ht=\paperheight,center,sep=2em]{palette primary}
    \usebeamerfont{title}Thank you.\\[0.5em]
    \usebeamerfont{subtitle}\small Generated from \texttt{hcw-brand.yaml} via \texttt{/new-theme}.
  \end{beamercolorbox}
\end{frame}

\end{document}
